\section*{अध्याय \ref{ch:g7:geological-time} --- पृथ्वी के युग}

\begin{solution}{exo:g7:geological-time:1}
(क) \omterm{def:g7:geological-time:unconformity}{विषमविन्यास}; (ख) अंतर्वेशन; (ग) अध्यारोपण; (घ) \omterm{def:g7:sediments-and-fossils:fossil}{जीवाश्म}
अनुक्रम; (ङ) अनुप्रस्थ काट।
\end{solution}

\begin{solution}{exo:g7:geological-time:2}
(क) लाल बलुआ पत्थर, सबसे नीचे, सबसे पुराना है; पीला बलुआ पत्थर, सबसे ऊपर, सबसे
नया है। (ख) \omterm{def:g7:geological-time:superposition}{अध्यारोपण का सिद्धांत}। (ग) बनने के बाद से परतें उलटी नहीं
गई हैं।
\end{solution}

\begin{solution}{exo:g7:geological-time:3}
(क) पुराजीवी, मध्यजीवी, नवजीवी। (ख) नवजीवी। (ग) मध्यजीवी। (घ) ट्राइऐसिक के ठीक
बाद जुरैसिक आता है; पैलियोजीन के ठीक पहले क्रिटेशियस
आता है।
\end{solution}

\begin{solution}{exo:g7:geological-time:4}
(क) ग़लत: \omterm{def:g7:geological-time:relative-dating}{सापेक्ष काल-निर्धारण} \omterm{def:g3:stones-pebbles-sand:rock}{चट्टानों} को केवल क्रम में रखता है; वर्षों में दी गई
आयु \omterm{def:g7:geological-time:relative-dating}{निरपेक्ष आयु} है। (ख) सही। (ग) ग़लत: \omterm{def:g3:stones-pebbles-sand:pebble}{कंकड़} उस कांग्लोमरेट से पुराने हैं जिसमें
वे हैं (\omterm{def:g7:geological-time:inclusion}{अंतर्वेशन का सिद्धांत})। (घ) सही (\omterm{def:g7:geological-time:fossil-succession}{जीवाश्म अनुक्रम का
सिद्धांत})।
\end{solution}

\begin{solution}{exo:g7:geological-time:5}
(क) नया: वह परत 3 को काटता है, और जो \omterm{def:g3:stones-pebbles-sand:rock}{चट्टान} दूसरी को काटती है वह नई होती है
(अनुप्रस्थ काट)। (ख) पुराना: परत 4 उसके ऊपर है और कटी नहीं है। (ग) परत 1,
परत 2, परत 3, ग्रेनाइट, परत 4। (घ) उसके ऊपर की \omterm{def:g3:stones-pebbles-sand:rock}{चट्टानें} \omterm{def:g5:water-shapes-the-land:erosion}{अपरदन} से तब तक घिसती रहीं
जब तक ग्रेनाइट सतह तक नहीं पहुँच गया; फिर परत 4 इसी \omterm{def:g5:water-shapes-the-land:erosion}{अपरदन} सतह पर बिछी, जो एक
\omterm{def:g7:geological-time:unconformity}{विषमविन्यास} है।
\end{solution}

\begin{solution}{exo:g7:geological-time:6}
(क) $538.8 - 251.9 = 286.9$ मिलियन वर्ष। (ख) $251.9 - 66 = 185.9$
मिलियन वर्ष। (ग) 66 मिलियन वर्ष। (घ) पुराजीवी सबसे लंबा है, नवजीवी सबसे छोटा
(अब तक)।
\end{solution}

\begin{solution}{exo:g7:geological-time:7}
(क) $4540 - 538.8 = 4001.2$ मिलियन वर्ष। (ख) $4001.2 \div 4540
\approx 0.881$:
लगभग \textbf{88~\%}। (ग) पुराजीवी से पहले सीपियों वाले जानवर अभी आम नहीं थे, और
नरम जीव लगभग कभी \omterm{def:g7:sediments-and-fossils:fossil}{जीवाश्म} नहीं बनते; सीपियाँ और हड्डियाँ, जो अच्छी तरह \omterm{def:g7:sediments-and-fossils:fossil}{जीवाश्म}
बनती हैं, केवल पुराजीवी से आगे ही आम
हैं।
\end{solution}

\begin{solution}{exo:g7:geological-time:8}
(क) $4540 \div 100 = 45.4$ मिलियन वर्ष प्रति मीटर। (ख) $538.8 \div 45.4
\approx 11.9$:
अंतिम रेखा से लगभग \qty{11.9}{m} पहले। (ग) $66 \div
45.4 \approx 1.45$: लगभग
\qty{1.45}{m}। (घ) $0.3 \div 45.4 \approx
0.0066$~m, यानी लगभग \qty{0.66}{cm},
एक सेंटीमीटर से भी कम: धावक का ध्यान शायद ही उस पर
जाए।
\end{solution}

\begin{solution}{exo:g7:geological-time:9}
(क) सिल्यूरियन पहले आया। हाँ: पुरानी \omterm{def:g3:stones-pebbles-sand:rock}{चट्टान} नई \omterm{def:g3:stones-pebbles-sand:rock}{चट्टान} के नीचे है। (ख) धूसर परतें
समतल बिछती हैं; वे धकेलकर खड़ी की जाती हैं; वे घिस जाती हैं; लाल \omterm{def:g3:stones-pebbles-sand:sand}{रेत} उनके ऊपर बिछती
है। (ग) यह बहुत अलग-अलग आयु की \omterm{def:g3:stones-pebbles-sand:rock}{चट्टानों} के बीच की एक \omterm{def:g5:water-shapes-the-land:erosion}{अपरदन} सतह है: जिस लंबे समय
में धूसर परतें ऊपर उठीं और घिसीं, उसने कोई परत पीछे नहीं
छोड़ी।
\end{solution}

\begin{solution}{exo:g7:geological-time:10}
(क) अमोनाइट वाले दोनों शेल एक ही आयु के हैं, और समुद्री अर्चिन वाले दोनों चूना
पत्थर भी: एक ही \omterm{def:g7:sediments-and-fossils:fossil}{जीवाश्म}, एक ही आयु (\omterm{def:g7:geological-time:fossil-succession}{जीवाश्म अनुक्रम का सिद्धांत})। (ख) सबसे
पुरानी परत \omterm{def:g3:stones-pebbles-sand:rock}{चट्टान} A का ट्राइलोबाइट वाला बलुआ पत्थर है (वह अमोनाइट वाले शेल के नीचे
है, जो \omterm{def:g3:stones-pebbles-sand:rock}{चट्टान} B की सबसे नीचे की परत जितना पुराना है)। सबसे नई परत \omterm{def:g3:stones-pebbles-sand:rock}{चट्टान} B की
पत्तियों वाली चिकनी मिट्टी है। (ग) नीचे से ऊपर: ट्राइलोबाइट वाला बलुआ पत्थर;
अमोनाइट वाला शेल; समुद्री अर्चिन वाला चूना पत्थर; पत्तियों वाली चिकनी
मिट्टी।
\end{solution}

\begin{solution}{exo:g7:geological-time:11}
(क) शिस्ट। वह कांग्लोमरेट के नीचे है (अध्यारोपण), और कांग्लोमरेट के \omterm{def:g3:stones-pebbles-sand:pebble}{कंकड़} शिस्ट के
टुकड़े हैं (अंतर्वेशन); \omterm{def:g5:water-shapes-the-land:erosion}{अपरदन} सतह भी शिस्ट को काटती है। (ख) बाद में: बलुआ पत्थर
समतल बिछा था, इसलिए वह बाद में झुका। (ग) पहले: चिकनी मिट्टी \omterm{def:g7:earthquakes:fault}{भ्रंश} को बिना टूटे
ढकती है। (घ) शिस्ट बनता है; वह घिसता है (लहरदार सतह); कांग्लोमरेट, बलुआ पत्थर और
चूना पत्थर, इसी क्रम में, बिछते हैं; परतें झुकती हैं, और \omterm{def:g7:earthquakes:fault}{भ्रंश} उन्हें तोड़ता है ---
अकेला अनुप्रस्थ काट यह नहीं बताता कि इन दोनों में से पहले क्या हुआ; झुकी परतें और
\omterm{def:g7:earthquakes:fault}{भ्रंश} घिसकर समतल हो जाते हैं; चिकनी मिट्टी ऊपर बिछती है। (ङ) दो: शिस्ट और
कांग्लोमरेट के बीच की लहरदार सतह, और झुकी परतों और समतल चिकनी मिट्टी के बीच
की सतह।
\end{solution}

\begin{solution}{exo:g7:geological-time:12}
(क) 365 दिनों का आधा 182.5 दिन है। जनवरी से जून तक
$31+28+31+30+31+30 = 181$ दिन होते हैं, इसलिए इतिहास जुलाई के 1.5 दिन बीतने पर आधा
होता है: \textbf{2~जुलाई को दोपहर 12 बजे}। (ख) $4540 \div 2 = 2270$ मिलियन
वर्ष। (ग) $80 \div \num{4540000000} \times \num{31536000} \approx 0.56$:
लगभग \textbf{0.6~सेकंड}।
\end{solution}

\begin{solution}{pb:g7:geological-time:1}
\textbf{1.} \num{4540} मिलियन वर्ष।

\textbf{2.} $4540 \div 365 \approx 12.4$ मिलियन वर्ष।

\textbf{3.} $\num{4540000000} \div (365 \times 24) = \num{4540000000}
\div 8760 \approx \num{518000}$
वर्ष (पूर्णांकित \num{12.4} मिलियन के साथ,
$\num{12400000} \div 24 \approx \num{517000}$)।

\textbf{4.} $\num{4540000000} \div \num{525600} \approx 8638$: लगभग
\num{8600} वर्ष।

\textbf{5.} $4540 - 4030 = 510$ मिलियन वर्ष, पृथ्वी बनने के बाद;
$510 \div 12.4 \approx 41$ दिन। जनवरी में 31 दिन हैं, इसलिए \textbf{11~फ़रवरी}
की शुरुआत में 41 दिन बीत चुके होते हैं।

\textbf{6.} $4540 - 3500 = 1040$ मिलियन वर्ष; $1040 \div 12.4 \approx
83.9$
दिन (बिना पूर्णांकित $12.44$ के साथ 83.6)। जनवरी और फ़रवरी में 59 दिन होते हैं:
मार्च के 24.6 से 24.9 दिन बीतने पर, \textbf{25~मार्च} को।

\textbf{7.} $66 \div 12.4 \approx 5.3$ दिन वर्ष के अंत से पहले, यानी
1~जनवरी के $365 - 5.3 = 359.7$ दिन बाद। पहले ग्यारह महीनों में 334 दिन होते हैं:
दिसंबर के 25.7 दिन बीतने पर, \textbf{26~दिसंबर} को।

\textbf{8.} $538.8 \div 4540 \approx 0.119$: वर्ष का \textbf{11.9~\%} बचता है,
यानी $0.119 \times 365 \approx 43$ दिन। $365 - 43 =
322$ दिन 1~जनवरी के बाद:
वह \textbf{नवंबर} में शुरू होता है (18~नवंबर
को)।

\textbf{9.} $\num{10000} \div 8600 \approx 1.2$ मिनट।

\textbf{10.} $5000 \div 8600 \approx 0.58$ मिनट, और $0.58 \times 60
\approx 35$:
आधी रात से लगभग \textbf{35 सेकंड} पहले।

\textbf{11.} $\num{300000} \div 8600 \approx 34.9$: आधी रात से लगभग \textbf{35
मिनट} पहले।

\textbf{12.} 31~दिसंबर की आधी रात से 35 मिनट पहले: हमारी जाति लगभग
\textbf{रात 11:25 बजे} (23:25) प्रकट होती है, वर्ष के आख़िरी घंटे में।
\end{solution}
